% GENERATED FILE. Edit canonical lessons, then run npm run generate:books.
% canonical-source-hash: fnv1a64:3d29a9d0717ca7e7
% canonical-lessons: TE-C469-samme, TE-C469-musiveta, TE-C469-antarayam, TE-C469-hadavidi, TE-C469-tarumaru, TE-R469-review-462-469

\chapter{A City at a Standstill}
\label{ch:te-a-city-at-a-standstill}
\hlchaptermodality{469}

\begin{chapteropening}
\textbf{By the end of this chapter:} \emph{I can say a strike, a closure, a disruption, a rush and a mix-up.}

Complete the last lesson of chapter 469: I can say a strike, a closure, a disruption, a rush and a mix-up.
\end{chapteropening}

\section[\texorpdfstring{\te{సమ్మె} (samme)}{సమ్మె (samme)}]{\te{సమ్మె} (samme) — a strike}

\label{lesson:TE-C469-samme}



\begin{quote}
Before the new one: say the Telugu for a slum, then the Telugu for to pile up.
\end{quote}

\subsection*{You'll want to know: \te{సమ్మె}}
\textbf{\te{సమ్మె}} — \emph{samme} — \textquotedblleft{}a strike\textquotedblright{}.

\begin{quote}
\te{బస్సు} \te{డ్రైవర్లు} \te{రెండు} \te{రోజులు} \te{సమ్మె} \te{చేశారు}.
\end{quote}

\begin{quote}
\emph{bassu ḍraivarlu reṇḍu rōjulu samme cēśāru.}
\end{quote}

\begin{quote}
The bus drivers went on strike for two days.
\end{quote}

\te{సమ్మె} \te{చేయు}, to strike; \te{సమ్మె} \te{విరమించు}, to call off a strike.

\begin{cousinweb}
\textbf{\te{సమ్మె}} (\emph{samme}) is a Telugu word for a strike whose origin is not settled.
\end{cousinweb}

\begin{grammarlens}[title={What you've built}]
One more word for getting around.
\end{grammarlens}

\subsection*{Your turn}
\begin{itemize}
\raggedright
  \item \emph{Say it:} \emph{samme}
  \item \emph{Say it:} \emph{samme}, once more
  \item \emph{Say it:} \emph{bassu ḍraivarlu reṇḍu rōjulu samme cēśāru.} — The bus drivers went on strike for two days.
  \item \emph{Recall:} say the Telugu for a slum, then the Telugu for to pile up, then say \emph{samme} again
\end{itemize}

\begin{tcolorbox}[breakable,colback=teal!4,colframe=teal!35!black,title={Before you move on}]
What does \te{సమ్మె} mean? (\textquotedblleft{}A strike\textquotedblright{}.) Say it once more.
\end{tcolorbox}
\section[\texorpdfstring{\te{మూసివేత} (mūsivēta)}{మూసివేత (mūsivēta)}]{\te{మూసివేత} (mūsivēta) — a closure; a shutdown}

\label{lesson:TE-C469-musiveta}



\begin{quote}
Before the new one: say the Telugu for to pile up, then the Telugu for a strike.
\end{quote}

\subsection*{You'll want to know: \te{మూసివేత}}
\textbf{\te{మూసివేత}} — \emph{mūsivēta} — \textquotedblleft{}a closure; a shutdown\textquotedblright{}.

\begin{quote}
\te{వంతెన} \te{మూసివేత} \te{వల్ల} \te{చుట్టూ} \te{తిరిగి} \te{వెళ్ళాం}.
\end{quote}

\begin{quote}
\emph{vantena mūsivēta valla cuṭṭū tirigi veḷḷāṁ.}
\end{quote}

\begin{quote}
Because the bridge was closed we went the long way round.
\end{quote}

\te{మూసి}, having shut, + -\te{వేత}, a putting: a closing that lasts.

\begin{cousinweb}
\textbf{\te{మూసివేత}} (\emph{mūsivēta}) is native Dravidian: \te{మూసి}, from \te{మూయు}, \textquotedblleft{}to cover, to shut\textquotedblright{} (DEDR 4915), and -\te{వేత}, the noun of \te{వేయు}, \textquotedblleft{}to put\textquotedblright{} (DEDR 5549).
\end{cousinweb}

\begin{grammarlens}[title={What you've built}]
One more word for getting around.
\end{grammarlens}

\subsection*{Your turn}
\begin{itemize}
\raggedright
  \item \emph{Say it:} \emph{mūsivēta}
  \item \emph{Say it:} \emph{mūsivēta}, once more
  \item \emph{Say it:} \emph{vantena mūsivēta valla cuṭṭū tirigi veḷḷāṁ.} — Because the bridge was closed we went the long way round.
  \item \emph{Recall:} say the Telugu for to pile up, then the Telugu for a strike, then say \emph{mūsivēta} again
\end{itemize}

\begin{tcolorbox}[breakable,colback=teal!4,colframe=teal!35!black,title={Before you move on}]
What does \te{మూసివేత} mean? (\textquotedblleft{}A closure; a shutdown\textquotedblright{}.) Say it once more.
\end{tcolorbox}
\section[\texorpdfstring{\te{అంతరాయం} (antarāyaṁ)}{అంతరాయం (antarāyaṁ)}]{\te{అంతరాయం} (antarāyaṁ) — a disruption; an interruption}

\label{lesson:TE-C469-antarayam}



\begin{quote}
Before the new one: say the Telugu for a strike, then the Telugu for a closure.
\end{quote}

\subsection*{You'll want to know: \te{అంతరాయం}}
\textbf{\te{అంతరాయం}} — \emph{antarāyaṁ} — \textquotedblleft{}a disruption; an interruption\textquotedblright{}.

\begin{quote}
\te{కరెంటు} \te{సరఫరాలో} \te{అంతరాయం} \te{కలిగింది}.
\end{quote}

\begin{quote}
\emph{kareṇṭu sarapharālō antarāyaṁ kaligindi.}
\end{quote}

\begin{quote}
There was a disruption in the power supply.
\end{quote}

\te{అంతరాయం} \te{కలుగు}, a disruption occurs. The TV says it when a programme breaks off.

\begin{cousinweb}
\textbf{\te{అంతరాయం}} (\emph{antarāyaṁ}) is borrowed from Sanskrit \emph{antarāya}, \textquotedblleft{}an intervention, an obstacle\textquotedblright{}, from \emph{antar-i}, \textquotedblleft{}to come between\textquotedblright{} (Monier-Williams).
\end{cousinweb}

\begin{grammarlens}[title={What you've built}]
One more word for getting around.
\end{grammarlens}

\subsection*{Your turn}
\begin{itemize}
\raggedright
  \item \emph{Say it:} \emph{antarāyaṁ}
  \item \emph{Say it:} \emph{antarāyaṁ}, once more
  \item \emph{Say it:} \emph{kareṇṭu sarapharālō antarāyaṁ kaligindi.} — There was a disruption in the power supply.
  \item \emph{Recall:} say the Telugu for a strike, then the Telugu for a closure, then say \emph{antarāyaṁ} again
\end{itemize}

\begin{tcolorbox}[breakable,colback=teal!4,colframe=teal!35!black,title={Before you move on}]
What does \te{అంతరాయం} mean? (\textquotedblleft{}A disruption; an interruption\textquotedblright{}.) Say it once more.
\end{tcolorbox}
\section[\texorpdfstring{\te{హడావిడి} (haḍāviḍi)}{హడావిడి (haḍāviḍi)}]{\te{హడావిడి} (haḍāviḍi) — a rush; a fuss}

\label{lesson:TE-C469-hadavidi}



\begin{quote}
Before the new one: say the Telugu for a closure, then the Telugu for a disruption.
\end{quote}

\subsection*{You'll want to know: \te{హడావిడి}}
\textbf{\te{హడావిడి}} — \emph{haḍāviḍi} — \textquotedblleft{}a rush; a fuss\textquotedblright{}.

\begin{quote}
\te{ఉదయం} \te{స్టేషన్లో} \te{చాలా} \te{హడావిడి} \te{ఉంది}.
\end{quote}

\begin{quote}
\emph{udayaṁ sṭēṣanlō cālā haḍāviḍi undi.}
\end{quote}

\begin{quote}
In the morning there is a great rush at the station.
\end{quote}

\te{హడావిడిగా}, in a rush. Also the fuss before a wedding or a festival.

\begin{cousinweb}
\textbf{\te{హడావిడి}} (\emph{haḍāviḍi}) is an everyday Telugu word for a bustle whose origin is not settled.
\end{cousinweb}

\begin{grammarlens}[title={What you've built}]
One more word for getting around.
\end{grammarlens}

\subsection*{Your turn}
\begin{itemize}
\raggedright
  \item \emph{Say it:} \emph{haḍāviḍi}
  \item \emph{Say it:} \emph{haḍāviḍi}, once more
  \item \emph{Say it:} \emph{udayaṁ sṭēṣanlō cālā haḍāviḍi undi.} — In the morning there is a great rush at the station.
  \item \emph{Recall:} say the Telugu for a closure, then the Telugu for a disruption, then say \emph{haḍāviḍi} again
\end{itemize}

\begin{tcolorbox}[breakable,colback=teal!4,colframe=teal!35!black,title={Before you move on}]
What does \te{హడావిడి} mean? (\textquotedblleft{}A rush; a fuss\textquotedblright{}.) Say it once more.
\end{tcolorbox}
\section[\texorpdfstring{\te{తారుమారు} (tārumāru)}{తారుమారు (tārumāru)}]{\te{తారుమారు} (tārumāru) — a mix-up; disorder}

\label{lesson:TE-C469-tarumaru}



\begin{quote}
Before the new one: say the Telugu for a disruption, then the Telugu for a rush.
\end{quote}

\subsection*{You'll want to know: \te{తారుమారు}}
\textbf{\te{తారుమారు}} — \emph{tārumāru} — \textquotedblleft{}a mix-up; disorder\textquotedblright{}.

\begin{quote}
\te{హోటలులో} \te{మా} \te{బుకింగు} \te{తారుమారు} \te{అయింది}.
\end{quote}

\begin{quote}
\emph{hōṭalulō mā bukiṅgu tārumāru ayindi.}
\end{quote}

\begin{quote}
Our booking at the hotel got mixed up.
\end{quote}

\te{తారుమారు} \te{అవు}, to get mixed up; \te{తారుమారు} \te{చేయు}, to turn upside down.

\begin{cousinweb}
\textbf{\te{తారుమారు}} (\emph{tārumāru}) is native Dravidian: DEDR 3194 lists Telugu \emph{tāṟumāṟu}, \textquotedblleft{}irregularity, disorder\textquotedblright{}, with Tamil \emph{tāṟumāṟu}, \textquotedblleft{}confusion\textquotedblright{}.
\end{cousinweb}

\begin{grammarlens}[title={What you've built}]
One more word for getting around.
\end{grammarlens}

\subsection*{Your turn}
\begin{itemize}
\raggedright
  \item \emph{Say it:} \emph{tārumāru}
  \item \emph{Say it:} \emph{tārumāru}, once more
  \item \emph{Say it:} \emph{hōṭalulō mā bukiṅgu tārumāru ayindi.} — Our booking at the hotel got mixed up.
  \item \emph{Recall:} say the Telugu for a disruption, then the Telugu for a rush, then say \emph{tārumāru} again
\end{itemize}

\begin{tcolorbox}[breakable,colback=teal!4,colframe=teal!35!black,title={Before you move on}]
What does \te{తారుమారు} mean? (\textquotedblleft{}A mix-up; disorder\textquotedblright{}.) Say it once more.
\end{tcolorbox}
\section[(dialogue)]{Words from chapters 462–469 — a review}

\label{lesson:TE-R469-review-462-469}



\begin{quote}
Say the Telugu for a counter and a branch office.
\end{quote}

\subsection*{Your turn}
Say these aloud:

\begin{itemize}
\raggedright
  \item a counter, a branch office, to fill in, to issue, to sanction
  \item an embassy, an ambassador, to allot, to extend, to renew
  \item to provide, to gather together, to organise, a supply, to collect payment
  \item a municipality, population, a resident, a local person, a stranger
  \item a town square, a public garden, a lawn, a street lamp, surroundings
  \item a storey, a basement, a compound wall, spacious, deserted
  \item sewage, sanitation, a stench, a slum, to pile up
  \item a strike, a closure, a disruption, a rush, a mix-up
\end{itemize}

\begin{tcolorbox}[breakable,colback=teal!4,colframe=teal!35!black,title={Before you move on}]
A counter? (\textbf{\te{కౌంటరు}}, \emph{kauṇṭaru}.) A branch office? (\textbf{\te{శాఖ}}, \emph{śākha}.) To fill in? (\textbf{\te{పూరించు}}, \emph{pūriñcu}.) To issue? (\textbf{\te{జారీ} \te{చేయు}}, \emph{jārī cēyu}.) To sanction? (\textbf{\te{మంజూరు} \te{చేయు}}, \emph{mañjūru cēyu}.) An embassy? (\textbf{\te{రాయబార} \te{కార్యాలయం}}, \emph{rāyabāra kāryālayaṁ}.) An ambassador? (\textbf{\te{రాయబారి}}, \emph{rāyabāri}.) To allot? (\textbf{\te{కేటాయించు}}, \emph{kēṭāyiñcu}.) To extend? (\textbf{\te{పొడిగించు}}, \emph{poḍigiñcu}.) To renew? (\textbf{\te{పునరుద్ధరించు}}, \emph{punaruddhariñcu}.) To provide? (\textbf{\te{అందించు}}, \emph{andiñcu}.) To gather together? (\textbf{\te{సమకూర్చు}}, \emph{samakūrcu}.) To organise? (\textbf{\te{ఏర్పాటు} \te{చేయు}}, \emph{ērpāṭu cēyu}.) A supply? (\textbf{\te{సరఫరా}}, \emph{sarapharā}.) To collect payment? (\textbf{\te{వసూలు} \te{చేయు}}, \emph{vasūlu cēyu}.) A municipality? (\textbf{\te{పురపాలక} \te{సంఘం}}, \emph{purapālaka saṅghaṁ}.) Population? (\textbf{\te{జనాభా}}, \emph{janābhā}.) A resident? (\textbf{\te{నివాసి}}, \emph{nivāsi}.) A local person? (\textbf{\te{స్థానికుడు}}, \emph{sthānikuḍu}.) A stranger? (\textbf{\te{అపరిచితుడు}}, \emph{aparicituḍu}.) A town square? (\textbf{\te{చౌరస్తా}}, \emph{caurastā}.) A public garden? (\textbf{\te{ఉద్యానవనం}}, \emph{udyānavanaṁ}.) A lawn? (\textbf{\te{పచ్చిక}}, \emph{paccika}.) A street lamp? (\textbf{\te{వీధిదీపం}}, \emph{vīdhidīpaṁ}.) Surroundings? (\textbf{\te{పరిసరాలు}}, \emph{parisarālu}.) A storey? (\textbf{\te{అంతస్తు}}, \emph{antastu}.) A basement? (\textbf{\te{నేలమాళిగ}}, \emph{nēlamāḷiga}.) A compound wall? (\textbf{\te{ప్రహరీ}}, \emph{praharī}.) Spacious? (\textbf{\te{విశాలమైన}}, \emph{viśālamaina}.) Deserted? (\textbf{\te{నిర్మానుష్యమైన}}, \emph{nirmānuṣyamaina}.) Sewage? (\textbf{\te{మురుగు}}, \emph{murugu}.) Sanitation? (\textbf{\te{పారిశుద్ధ్యం}}, \emph{pāriśuddhyaṁ}.) A stench? (\textbf{\te{దుర్వాసన}}, \emph{durvāsana}.) A slum? (\textbf{\te{మురికివాడ}}, \emph{murikivāḍa}.) To pile up? (\textbf{\te{పేరుకుపోవు}}, \emph{pērukupōvu}.) A strike? (\textbf{\te{సమ్మె}}, \emph{samme}.) A closure? (\textbf{\te{మూసివేత}}, \emph{mūsivēta}.) A disruption? (\textbf{\te{అంతరాయం}}, \emph{antarāyaṁ}.) A rush? (\textbf{\te{హడావిడి}}, \emph{haḍāviḍi}.) A mix-up? (\textbf{\te{తారుమారు}}, \emph{tārumāru}.)
\end{tcolorbox}
